\documentclass[11pt,a4paper]{article}

\usepackage[margin=2.3cm]{geometry}
\usepackage[T1]{fontenc}
\usepackage[utf8]{inputenc}
\usepackage{fourier}
\usepackage{microtype}
\usepackage{enumitem}
\usepackage[dvipsnames]{xcolor}
\usepackage[most]{tcolorbox}
\usepackage[colorlinks=true,linkcolor=black,urlcolor=NavyBlue]{hyperref}

\setlength{\parindent}{0pt}
\setlength{\parskip}{0.5em}

\newcommand{\code}[1]{\texttt{#1}}

\title{\vspace{-1.5em}HW1:  README\\
       \large ML for Finance (Prof. Arratia), BSE 2026}
\author{Marco Canal \and Michael Duarte Gonçalves \and Olivia Peter}
\date{15 May 2026}

\begin{document}
\maketitle

\section*{Project layout}

The full report is \code{solution.pdf} (from \code{solution.tex}). Code lives
in \code{scripts/}. The committed \code{data/} folder contains the frozen
\code{goyal-welch2022Monthly.csv} panel. We keep it under version control
for reproducibility, so that Exercises 2 and 3 produce the same numbers
without an internet round-trip. The \code{output/} folders (\code{ex1/},
\code{ex2/}, \code{ex3/}) are created on demand by the scripts and hold
every CSV, PDF figure and TeX fragment cited in the report.

\section*{The four scripts}

\textbf{\code{exercise\_01.py}} covers Exercise 1. It pulls daily prices
for ten assets from Yahoo Finance over our team epoch (2024-01-01 to
2025-12-31), aligns crypto to the NYSE calendar, computes log returns and
RiskMetrics EWMA volatility ($\lambda=0.94$) at weekly and monthly
frequency, then for every ordered pair selects the AIC-optimal VAR lag in
$[1,5]$ and runs a Granger F-test at that lag. It produces the four
causality tables and the four network plots in \code{output/ex1/}.

\textbf{\code{exercise\_02.ipynb}} covers Exercise 2. It trains the NN,
LSTM and GP grids on the Goyal-Welch panel (\code{dp}, \code{svar},
\code{bm}, plus lags of the target) across the three target specifications,
together with the drop-each-indicator sensitivity ablation. It writes to
\code{output/ex2/}.

\textbf{\code{exercise\_03.py}} covers Exercise 3. It fits the plain and
Yeo-Johnson ARMA-GARCH baselines on the same chronological split as
Exercise 2 and scores them on the same 285-observation test block, so the
head-to-head table is apples-to-apples. It writes to \code{output/ex3/}.

\section*{Running the pipeline end-to-end}

From a fresh checkout, with Python 3.11 and a TeX Live install:

\begin{enumerate}[leftmargin=*,itemsep=0.1em]
\item \code{pip install -r requirements.txt}
\item \code{python scripts/exercise\_01.py} (creates \code{output/ex1/})
\item Run \code{scripts/exercise\_02.ipynb} (creates \code{output/ex2/}, see GPU note below)
\item \code{python scripts/exercise\_03.py} (creates \code{output/ex3/})
\end{enumerate}

Measured wall-times. Each script logs its own wall time at the
end of the run. The numbers below are from our last full run on one of
the authors' laptops: an ASUS ROG Strix G18, Intel
i9-14900HX (32 logical cores, 32~GB RAM, Ubuntu 24.04) with an
NVIDIA RTX 4080 Laptop GPU (12~GB VRAM) used for the TensorFlow
steps. They are indicative only: timings vary with CPU generation, disk
speed, Yahoo Finance latency and the GPU on hand.

\vspace{12pt}

\begin{itemize}[leftmargin=*,itemsep=0.05em]
\item[] \code{exercise\_01.py}: 3.7~s.
\item[] \code{exercise\_02.ipynb}: 13.6~min end to end on the local
      RTX 4080. For reference, the local GPU is generally faster than
      Colab on these small NN and LSTM tensors. Per-cell wall times also
      print via the \code{\%\%time}.
\item[] \code{exercise\_03.py}: 7.0~s. (the \code{auto\_arima}
      stepwise search and the GARCH MLE on the percent-scale residuals
      dominate the runtime).
\end{itemize}

\begin{tcolorbox}[colback=NavyBlue!4,colframe=NavyBlue!70!black,
                  title=\textbf{Note: running Exercise 2 (the GPU step)},
                  fonttitle=\bfseries,boxrule=0.6pt,arc=2pt,
                  left=6pt,right=6pt,top=4pt,bottom=4pt]

The NN, LSTM and GP grids in \code{exercise\_02.ipynb} finish in a few
minutes on a GPU and would take several hours on CPU. The path-setup
cell auto-detects Colab vs.\ a local install.

\textbf{Option A. Google Colab.} Upload \code{HW1/} to Drive at
\code{MyDrive/ML4F\_HW1}, open the notebook from Drive in Colab,
\emph{Runtime} $\rightarrow$ \emph{Change runtime type} $\rightarrow$
select a T4 GPU, then \emph{Runtime} $\rightarrow$ \emph{Run all}.

\vspace{12pt}

\textbf{Option B. Local NVIDIA GPU.} The base \code{requirements.txt}
pins plain \code{tensorflow}, which is enough for CPU. For GPU
acceleration, layer the bundled CUDA wheels on top with
\code{pip install --upgrade tensorflow[and-cuda]}. That pulls in
\code{nvidia-cuda-runtime-cu12}, \code{nvidia-cublas-cu12},
\code{nvidia-cudnn-cu12} and friends, so TF self-contains its runtime
and no system toolkit is required.

\end{tcolorbox}

A small caveat on reproducibility: we fix \code{SEED=7} for NumPy,
TensorFlow and PyTorch, but GPU kernels and cuDNN convolutions are not
fully deterministic. NRMSE values can move in the third decimal place on
a different machine. The report quotes four decimal places only where the
ranking is robust to that level of jitter.

\end{document}
